\documentclass[10pt,a4paper]{amsart}
\usepackage[margin=19mm]{geometry}
\usepackage{amsmath,amssymb,mathtools}
\usepackage[scr=rsfs,cal=euler]{mathalfa}
\usepackage{xfrac}
\usepackage[unicode,hidelinks,bookmarks=false]{hyperref}
\newcommand{\ie}{i.e.\ }
\newcommand{\cf}{cf.\ }
\newcommand{\resp}{respectively\ }
\newcommand{\Subsection}{\subsection}
\providecommand{\nobreakdash}{\nobreak\hbox{-}}
\setlength{\emergencystretch}{2em}
\multlinegap0pt
\usepackage{amssymb}
\usepackage[shortlabels]{enumitem}
\usepackage{float}
\usepackage{xcolor}
\newtheorem{lemma}{Lemma}
\newtheorem{corollary}{Corollary}
\title{Matchings on Random Regular Hypergraphs}
\author{Zhongyang Li}
\date{Source: arXiv:2108.11003v3 (CC BY 4.0)}
\begin{document}
\maketitle

\noindent\emph{Source and licence.} Matchings on Random Regular Hypergraphs. Version arXiv:2108.11003v3 (CC BY 4.0), distributed by arXiv under the Creative Commons Attribution 4.0 International licence, http://creativecommons.org/licenses/by/4.0/, as recorded in the arXiv licence metadata for this exact version and shown on the abstract page https://arxiv.org/abs/2108.11003v3. Full licence notice: \texttt{licenses/arxiv-2108.11003-CC-BY-4.0.txt}.\par\smallskip

\noindent\textit{Editorial scope.} Counted unit: the article's lemma l51 (source0.tex lines 2380--2393). The article's complete original proof of it is delivered with it (source0.tex lines 2395--2552, one \texttt{proof} environment of 158 lines). No mathematics is written, repaired or re-derived: the statement and the proof are the article's own lines, in the article's own notation, and no retained line is substituted or re-flowed.  Retained uncounted as the source-local context the proof needs: lines 517--519; lines 742--762; lines 811--825; lines 1699--1704; lines 1794--1805; lines 1926--1935; lines 2069--2078; lines 2100--2106; lines 2169--2201; lines 2280--2290; lines 2328--2336.  Nothing was omitted from any retained span, so the package carries no omission record.  No retained line contains a citation, so the wrapper carries no bibliography and no bibliography entry.  No retained line contains a figure environment, an image-inclusion command or drawing code, so no image is delivered and the package declares no asset dependency.  Editorial formatting only, in addition: an AMS \texttt{amsart} preamble that restates the article's own declarations and macros used by the retained lines, namely the environments \texttt{lemma}, \texttt{corollary} on separate counters, together with the standard packages those lines need.  The retained environments print in this document as Lemma 1, Corollary 1 in order of appearance, whereas the article prints the same items with the numbering of its own full document.  The complete native source of the article as distributed in the arXiv e-print for this exact version is bundled unchanged as \texttt{sources/115297/source0.tex}.
\begin{equation}\label{eq1}
  (1-d\beta)^l=\beta(1-\beta)^{l-1}.
\end{equation}

\begin{lemma}\label{lm21}
The following hold.
\begin{enumerate}
\item $\Phi_{d,l}$ has a unique maximizer $\beta_*\in(0,1/d)$, and $\beta_*$ is the solution of \eqref{eq1}.  In particular,
\[
\max_{0<\beta<1/d}\Phi_{d,l}(\beta)=\Phi_{d,l}(\beta_*).
\]
\item If
\[
\frac{(l-1)(d-1)}{d}\ln\left(\frac{d-1}{d}\right)
-\frac1d\ln\left(\frac1d\right)
=:f_l\left(\frac1d\right)\geq 0,
\]
then $\Phi_{d,l}(\beta)>0$ for every $\beta\in(0,1/d)$.
\item If $f_l(1/d)<0$, then there is a unique $\beta_0\in(\beta_*,1/d)$ such that
\[
\Phi_{d,l}(\beta_0)=0.
\]
Moreover, $\Phi_{d,l}(\beta)>0$ for $\beta\in(0,\beta_0)$, and $\Phi_{d,l}(\beta)<0$ for $\beta\in(\beta_0,1/d)$.
\end{enumerate}
\end{lemma}

\begin{lemma}\label{lem:first-moment-local}
Let $K\subset(0,1/d)$ be compact.  Then
\[
\sup_{\substack{h\in\mathbb{N}:\ h/m\in K}}
\left|
\sqrt{m}\,e^{-m\Phi_{d,l}(h/m)}\mathbb{E}Z_h
-\frac{1}{\sqrt{2\pi(h/m)(1-dh/m)}}
\right|\longrightarrow 0.
\]
In particular, if $\beta\in(0,1/d)$ is fixed and $m\beta\in\mathbb{N}$, then
\[
\lim_{m\to\infty}\sqrt{m}\,e^{-m\Phi_{d,l}(\beta)}\mathbb{E}Z_{m\beta}
=\frac{1}{\sqrt{2\pi\beta(1-d\beta)}}.
\]
\end{lemma}

\begin{equation}
\lim_{m\to\infty}
\frac{\mathbb E Z_{m\beta}^2}{(\mathbb E Z_{m\beta})^2}
=\frac{1-\beta}{\sqrt{\beta^2(d+l-dl)-2\beta+1}}.
\label{cm}
\end{equation}

\begin{corollary}[Compact-uniform second-moment ratio]\label{cor:uniform-second-ratio-L1}
Let $K\subset(0,1/d)$ be compact.  Suppose that either $l=2$, or $l\ge3$ and
$K\subset(0,(dl-d-l+2)^{-1}]$ with $\Phi_{d,l}(\beta)>0$ for every $\beta\in K$.
Then the convergence in \eqref{cm} is uniform for $\beta\in K$ along admissible
subsequences $m\beta\in\mathbb N$.  Equivalently, if $h_m/m\in K$ and
$h_m/m\to\beta\in K$, then
\[
  \frac{\mathbb E Z_{h_m}^2}{(\mathbb E Z_{h_m})^2}
  \longrightarrow
  \frac{1-\beta}{\sqrt{\beta^2(d+l-dl)-2\beta+1}}.
\]
\end{corollary}

\begin{lemma}\label{l30}
For fixed \(k\),
\begin{equation}\label{eq:cond-cycle-mean}
  \mathbb E(C_k\mid A_{M_0})
  \longrightarrow
  \mu_{\beta,k}
  :=\frac{(d-1)^k(l-1)^k}{2k}
    \left(1+\frac{(-1)^k\beta^k}{(1-\beta)^k}\right).
\end{equation}
\end{lemma}

\begin{lemma}\label{l33}
For each fixed \(b\in\mathbb N\) the following hold as \(m\to\infty\).
\begin{enumerate}
\item The vector \((C_1,\ldots,C_b)\) converges in distribution to a vector of
independent Poisson random variables with means \(\lambda_1,\ldots,\lambda_b\).
\item Conditional on \(A_{M_0}\), the vector \((C_1,\ldots,C_b)\) converges in
distribution to a vector of independent Poisson random variables with means
\(\mu_{\beta,1},\ldots,\mu_{\beta,b}\).
\end{enumerate}
\end{lemma}

\begin{equation}\label{ebc}
  E_{\beta,\mathbf c}
  =\left(\prod_{k=1}^b
      e^{\lambda_k-\mu_{\beta,k}}
      \left(\frac{\mu_{\beta,k}}{\lambda_k}\right)^{c_k}\right)
    \mathbb E Z_{m\beta}\,(1+o(1)).
\end{equation}

\begin{lemma}\label{l35}
Assume \(r_\beta<1\).  Fix \(b\in\mathbb N\) and \(W>2\).  Put
\[
  \gamma_\beta:=1-q_\beta,
  \qquad A_\beta:=1+q_\beta,
  \qquad \lambda_{*,b}:=\min_{1\le k\le b}\lambda_k,
\]
and
\begin{align*}
\Theta_{\beta,b,W}(y)
:=1&-b\exp\left\{-\frac{\gamma_\beta^2y^2}{8A_\beta^4}\right\}
    -b\exp\left\{-\frac{\gamma_\beta^2y^2}{16A_\beta^4}\right\} \\
&-b\exp\left\{-\frac{|\phi(1/2)|\gamma_\beta^2y^2}{4A_\beta^4W^2}\right\}
    -b\exp\left\{-(\log(1+W)-1)
       \frac{\gamma_\beta^2y\sqrt{\lambda_{*,b}}}{2A_\beta^2}\right\}.
\end{align*}
Then, for all sufficiently large \(y\),
\begin{align*}
\liminf_{m\to\infty}
\frac{\sum_{\mathbf c\in S_b(y)}p_{\mathbf c}E_{\beta,\mathbf c}^2}
     {(\mathbb E Z_{m\beta})^2}
\ge
\Theta_{\beta,b,W}(y)
\left(1-r_\beta^b\right)
\frac1{\sqrt{1-r_\beta}},
\end{align*}
where \(p_{\mathbf c}=\mathbb P((C_1,\ldots,C_b)=\mathbf c)\).
Equivalently,
\[
  \frac1{\sqrt{1-r_\beta}}
  =\frac{1-\beta}{\sqrt{1-2\beta-(dl-d-l)\beta^2}}.
\]
\end{lemma}

\begin{lemma}\label{l38}
Assume \(r_\beta<1\), equivalently
\[
  \beta<\frac1{1+\sqrt{(d-1)(l-1)}}.
\]
There exist constants \(A,B>0\), depending only on \(d,l,\beta\), such that for
every fixed \(b\), every \(\mathbf c\in S_b(y)\), and all sufficiently large \(m\),
\[
  E_{\beta,\mathbf c}\ge e^{-(A+By)}\,\mathbb E Z_{m\beta}.
\]
\end{lemma}

\begin{lemma}[Compact-uniform cycle estimates]\label{lem:cycle-uniformity}
Fix $b\in\mathbb N$ and fixed cycle lengths and multiplicities.  Let
$K\subset(0,1/d)$ be compact.  Then the limits in Lemmas~\ref{l30}--\ref{l33}
and the fixed-$b$ asymptotic \eqref{ebc} hold uniformly for
$\beta\in K$ along admissible subsequences $m\beta\in\mathbb N$.  If, in addition,
$r_\beta$ is bounded away from $1$ on $K$, then the estimates in Lemmas~\ref{l35}
and~\ref{l38} hold uniformly for $\beta\in K$ after allowing the constants in
Lemma~\ref{l38} to depend on $K$.
\end{lemma}

\begin{lemma}\label{l51}
Let $h_m\in\mathbb N$, put $\beta_m=h_m/m$, and assume that
$\beta_m\to\beta\in(0,1/d)$.  Suppose that one of the following alternatives holds:
\begin{enumerate}
\item $l=2$; or
\item $l\ge3$ and $\beta_m\le L_1$ for all sufficiently large $m$.
\end{enumerate}
Then
\begin{equation}\label{cfe}
  \frac1m\log Z_{h_m}\longrightarrow \Phi_{d,l}(\beta)
\end{equation}
in probability.  In particular, if $m\beta\in\mathbb N$ along a subsequence, the
same conclusion holds with $h_m=m\beta$.
\end{lemma}

\begin{proof}
We first check that $\Phi_{d,l}(\beta)>0$.  If $l=2$, this follows from
Lemma \ref{lm21}(2), since $f_2(1/d)\ge0$.  If $l\ge3$, then
$L_1=1/(dl-d-l+2)$.  Put
\[
B:=dl-d-l+2=(d-1)(l-1)+1.
\]
For $0<u<1$ we have $(1-u)\log(1-u)\ge -u$ and
$-(1-u)\log(1-u)\ge u(1-u)$.  Hence, for every $0<\beta'\le 1/B$,
\begin{align*}
\Phi_{d,l}(\beta')
&\ge \beta'\log\frac1{\beta'}-(l-1)\beta'
   +l\beta'(1-d\beta') \\
&=\beta'\left(\log\frac1{\beta'}+1-ld\beta'\right)
\ge \beta'\left(\log B+1-\frac{ld}{B}\right)>0.
\end{align*}
The last inequality uses $B\ge3$ and $ld/B\le2$, which is equivalent to
$(d-2)(l-2)\ge0$.  Since $\beta_m\to\beta$ and $\beta_m\le L_1$ eventually in
case $l\ge3$, this proves the desired positivity.

By the compact-uniform first moment estimate in Lemma \ref{lem:first-moment-local},
\begin{equation}\label{eq:section5-first-moment-log}
  \frac1m\log \mathbb E Z_{h_m}\longrightarrow \Phi_{d,l}(\beta).
\end{equation}
It remains to prove
\begin{equation}\label{lmb}
  \frac1m\log\frac{Z_{h_m}}{\mathbb E Z_{h_m}}\longrightarrow0
  \quad\text{in probability}.
\end{equation}
The assumptions also imply $r_\beta<1$.  If $l=2$, then for $d=2$ we have
$r_\beta=q_\beta^2<1$, while for $d\ge3$ the inequality $\beta<1/d$ gives
$q_\beta<1/(d-1)$ and hence $r_\beta=(d-1)q_\beta^2<1$.  If $l\ge3$, then
$(d-1)(l-1)>1$ and the bound $\beta_m\le L_1$ implies
$q_\beta\le 1/((d-1)(l-1))$, whence
$r_\beta\le 1/((d-1)(l-1))<1$.

Write
\[
  X_m:=Z_{h_m},\qquad \mathcal F_{m,b}:=\sigma(C_1,\ldots,C_b),\qquad
  X_{m,b}:=\mathbb E(X_m\mid\mathcal F_{m,b}).
\]
All cycle estimates below are applied with $\beta=\beta_m$; the compact-uniform
form needed for this varying-density use is supplied by Lemma~\ref{lem:cycle-uniformity}.

Fix $b\ge1$.  For every fixed $\mathbf c=(c_1,\ldots,c_b)\in\mathbb N^b$, the
asymptotic formula \eqref{ebc}, applied with $\beta_m$ in place of $\beta$, gives
\begin{equation}\label{eq:section5-Ebc-limit}
  \frac{E_{\beta_m,\mathbf c}}{\mathbb E Z_{h_m}}
  \longrightarrow
  \prod_{k=1}^b e^{-\lambda_k\delta_k}(1+\delta_k)^{c_k},
  \qquad
  \delta_k:=(-1)^kq_\beta^k .
\end{equation}
The convergence is uniform over finite sets of $\mathbf c$'s by
Lemma~\ref{lem:cycle-uniformity}.  Lemma
\ref{l33}(1) gives
\[
  (C_1,\ldots,C_b)\Rightarrow (P_1,\ldots,P_b),
\]
where $P_1,\ldots,P_b$ are independent Poisson random variables with means
$\lambda_1,\ldots,\lambda_b$.  Combining this convergence with
\eqref{eq:section5-Ebc-limit} and tightness of $(C_1,\ldots,C_b)$ yields
\begin{equation}\label{eq:section5-Xmb-limit}
  \frac{X_{m,b}}{\mathbb E X_m}\Rightarrow
  Y_b:=\prod_{k=1}^b e^{-\lambda_k\delta_k}(1+\delta_k)^{P_k}.
\end{equation}
Since $q_\beta<1$, all factors in $Y_b$ are positive.

The sequence $Y_b$ converges almost surely to a finite strictly positive random
variable.  Indeed,
\begin{align*}
\log Y_b
&=\sum_{k=1}^b (P_k-\lambda_k)\log(1+\delta_k)
  +\sum_{k=1}^b \lambda_k\{\log(1+\delta_k)-\delta_k\}.
\end{align*}
Because
\begin{equation}\label{eq:section5-l2-summable}
  \sum_{k\ge1}\lambda_k\delta_k^2
  =\sum_{k\ge1}\frac{r_\beta^k}{2k}<\infty,
\end{equation}
the deterministic series is absolutely convergent, and the centred independent
series is convergent in $L^2$ and almost surely.  We denote the limit by
$Y_\infty$; then $Y_\infty\in(0,\infty)$ almost surely.

We next show that the residual $X_m-X_{m,b}$ is negligible on the exponential
scale.  Since $X_{m,b}=\mathbb E(X_m\mid\mathcal F_{m,b})$,
\begin{equation}\label{eq:section5-conditional-var-identity}
  \mathbb E(X_m-X_{m,b})^2
  =\mathbb E X_m^2-\mathbb E X_{m,b}^2.
\end{equation}
The hypotheses of the lemma imply the second moment asymptotic \eqref{cm}; for
$\beta_m\to\beta$ the convergence is uniform by
Corollary~\ref{cor:uniform-second-ratio-L1}.  Hence
\begin{equation}\label{eq:section5-second-ratio}
  \frac{\mathbb E X_m^2}{(\mathbb E X_m)^2}
  \longrightarrow
  R_\beta:=\frac1{\sqrt{1-r_\beta}}.
\end{equation}
On the other hand, fix any $W_0>2$.  Lemma \ref{l35}, applied with $W=W_0$
and with $\beta_m$ in place of $\beta$, and then using the compact-uniform form in
Lemma~\ref{lem:cycle-uniformity} together with $\beta_m\to\beta$,
gives, for every fixed $b$ and after letting its truncation parameter $y\to\infty$,
\begin{equation}\label{eq:section5-truncated-second-lower}
  \liminf_{m\to\infty}\frac{\mathbb E X_{m,b}^2}{(\mathbb E X_m)^2}
  \ge (1-r_\beta^b)R_\beta .
\end{equation}
Combining \eqref{eq:section5-conditional-var-identity},
\eqref{eq:section5-second-ratio}, and \eqref{eq:section5-truncated-second-lower},
we obtain
\begin{equation}\label{eq:section5-residual-variance}
  \limsup_{m\to\infty}
  \frac{\mathbb E(X_m-X_{m,b})^2}{(\mathbb E X_m)^2}
  \le R_\beta r_\beta^b .
\end{equation}

We now prove the two tails in \eqref{lmb}.  The upper tail follows immediately
from Markov's inequality:
\begin{equation}\label{eq:section5-upper-tail}
  \mathbb P\left\{X_m>e^{m\varepsilon}\mathbb E X_m\right\}
  \le e^{-m\varepsilon}.
\end{equation}
For the lower tail, let $\alpha>0$ be arbitrary.  Choose $a>0$ such that
\[
  \mathbb P\{Y_\infty\le 4a\}<\alpha/3.
\]
Then choose $b$ so large that
\begin{equation}\label{eq:section5-choose-b}
  \mathbb P\{|Y_b-Y_\infty|>a\}<\alpha/3,
  \qquad
  \frac{4R_\beta r_\beta^b}{a^2}<\alpha/3.
\end{equation}
For this fixed $b$, the Portmanteau theorem, \eqref{eq:section5-Xmb-limit}, and \eqref{eq:section5-choose-b} give
\begin{align}\label{eq:section5-Xmb-lower-tight}
  \limsup_{m\to\infty}
  \mathbb P\left\{X_{m,b}<2a\,\mathbb E X_m\right\}
  &\le \mathbb P\{Y_b\le 2a\} \notag\\
  &\le \mathbb P\{Y_\infty\le 4a\}+\mathbb P\{|Y_b-Y_\infty|>a\}\le 2\alpha/3.
\end{align}
If $m$ is large enough that $e^{-m\varepsilon}<a$, then
\begin{align*}
\mathbb P\left\{X_m<e^{-m\varepsilon}\mathbb E X_m\right\}
&\le
\mathbb P\left\{X_{m,b}<2a\,\mathbb E X_m\right\} \\
&\quad+
\mathbb P\left\{|X_m-X_{m,b}|>a\,\mathbb E X_m\right\}.
\end{align*}
By Chebyshev's inequality and \eqref{eq:section5-residual-variance}, the limsup of
the second probability is at most $R_\beta r_\beta^b/a^2<\alpha/12$.  Together with
\eqref{eq:section5-Xmb-lower-tight}, this gives
\[
  \limsup_{m\to\infty}
  \mathbb P\left\{X_m<e^{-m\varepsilon}\mathbb E X_m\right\}
  \le \alpha .
\]
Since $\alpha>0$ is arbitrary, the lower tail tends to zero.  Combining this with
\eqref{eq:section5-upper-tail} proves \eqref{lmb}.  Finally,
\eqref{eq:section5-first-moment-log} and \eqref{lmb} imply \eqref{cfe}.
\end{proof}

\end{document}
